\documentclass[10pt,a4paper]{amsart}
\usepackage[margin=19mm]{geometry}
\usepackage{amsmath,amssymb,mathtools}
\usepackage[scr=rsfs,cal=euler]{mathalfa}
\usepackage{xfrac}
\usepackage[unicode,hidelinks,bookmarks=false]{hyperref}
\newcommand{\ie}{i.e.\ }
\newcommand{\cf}{cf.\ }
\newcommand{\resp}{respectively\ }
\newcommand{\Subsection}{\subsection}
\providecommand{\nobreakdash}{\nobreak\hbox{-}}
\setlength{\emergencystretch}{2em}
\multlinegap0pt
\usepackage{amssymb}
\usepackage{algorithm}\usepackage{algpseudocode}
\newtheorem{proposition}{Proposition}
\newcommand{\At}{\mathscr{A}}
\newcommand{\Bv}{\mathbf{B}}
\newcommand{\Es}{\mathcal{E}}
\newcommand{\Hs}{\mathcal{H}}
\newcommand{\Nv}{\mathbf{N}}
\newcommand{\Rv}{\mathbf{R}}
\newcommand{\Ss}{\mathcal{S}}
\newcommand{\Vs}{\mathcal{V}}
\newcommand{\rank}{\text{rank}}
\title{\fontsize{22}{20}\selectfont Structural Controllability of Large-Scale Hypergraphs}
\author{Joshua Pickard, Xin Mao,, Can Chen}
\date{Source: arXiv:2603.19955v1 (CC BY 4.0)}
\begin{document}
\maketitle

\noindent\emph{Source and licence.} \fontsize{22}{20}\selectfont Structural Controllability of Large-Scale Hypergraphs. Version arXiv:2603.19955v1 (CC BY 4.0), distributed by arXiv under the Creative Commons Attribution 4.0 International licence, http://creativecommons.org/licenses/by/4.0/, as recorded in the arXiv licence metadata for this exact version and shown on the abstract page https://arxiv.org/abs/2603.19955v1. Full licence notice: \texttt{licenses/arxiv-2603.19955-CC-BY-4.0.txt}.\par\smallskip

\noindent\textit{Editorial scope.} Counted unit: the article's proposition polynomials (source0.tex lines 848--850). The article's complete original proof of it is delivered with it (source0.tex lines 853--935, one \texttt{proof} environment of 83 lines). No mathematics is written, repaired or re-derived: the statement and the proof are the article's own lines, in the article's own notation, and no retained line is substituted or re-flowed.  No further source-local context is needed: every cross-reference of every retained line resolves inside the two delivered spans.  Nothing was omitted from any retained span, so the package carries no omission record.  No retained line contains a citation, so the wrapper carries no bibliography and no bibliography entry.  No retained line contains a figure environment, an image-inclusion command or drawing code, so no image is delivered and the package declares no asset dependency.  Editorial formatting only, in addition: an AMS \texttt{amsart} preamble that restates the article's own declarations and macros used by the retained lines, namely the environments \texttt{proposition} on separate counters, together with the standard packages those lines need.  The retained environments print in this document as Proposition 1 in order of appearance, whereas the article prints the same items with the numbering of its own full document.  The complete native source of the article as distributed in the arXiv e-print for this exact version is bundled unchanged as \texttt{sources/196818/source0.tex}.
\begin{proposition}\label{thm: structural control of polynomials}
   The system $(\mathscr A,\textbf B)$ is structurally controllable if and only if the hypergraph $\Hs(\mathscr A,\textbf B)$ contains neither a hyperedge dilation nor an inaccessible node.
\end{proposition}

\begin{proof}
    We first show the sufficiency.
    %
    If $(\mathscr A,\textbf B)$ contains a nonaccessible node $v$, then there exists no walk beginning at the control nodes that will reach node $v$.
    %
    This indicates there will be a row of all zeros in the matrix $\textbf W_j$ for all possible steps $j$ in the walk.
    %
    There will then be a row of all zeros in $\textbf W.$
    %
    Since $\textbf W$ has the same sparsity structure as the nonlinear controllability matrix, then for any pair $(\tilde{\mathscr A},\tilde{\textbf B})$ the nonlinear controllability matrix will have a row of all zeros and be low rank.
    %
    Therefore, the pair $(\mathscr A,\textbf B)$ is not structurally controllable.
If $\Hs(\mathscr A,\textbf B)$ contains a dilation, the augmented matrix $\textbf R=\begin{bmatrix}
        \tilde{\textbf A}_{(1)}&\tilde{\textbf B}
    \end{bmatrix}$ may be partitioned as
    \begin{equation*}
        \textbf R=\begin{bmatrix}
            \tilde{\textbf A}_{(1)}&\tilde{\textbf B}    \end{bmatrix}=\begin{bmatrix}
            \textbf D_\mathscr A&\textbf D_{\textbf B}\\ \textbf N_\mathscr A&\textbf N_\textbf B
        \end{bmatrix},
    \end{equation*}
    where $\textbf D=\begin{bmatrix}
        \textbf D_\mathscr A&\textbf D_\textbf B
    \end{bmatrix}$ is a $d\times (n^{k-1}+m)$ matrix whose rows correspond to nodes in the dilation set $\Ss$, and $\textbf N=\begin{bmatrix}
        \textbf N_\mathscr A&\textbf N_\textbf B
    \end{bmatrix}$ is a $(n-d)\times (n^{k-1}+m)$ matrix whose rows correspond to nodes not found within the dilation.
    %
    Since $\Ss$ is a dilated set of nodes, there exists $d-1$ or fewer hyperedges whose heads are unique with respect to the elements in $\Ss.$
    %
    As the rows of $\textbf D$ are nodes and the columns are hyperedges, there exists at most $d-1$ unique column vectors in $\textbf D,$ so $\rank(\textbf D)<d.$
    %
    The matrix $\textbf W$ has the sparsity structure of
    \begin{equation*}
    \begin{split}
        \textbf W&=\begin{bmatrix}
            \textbf D_\textbf B 
        & \textbf D_\mathscr A\textbf W_0^{[k-1]}&\cdots&\textbf D_\mathscr A\textbf W_{n-1}^{[k-1]}\\
            &&\\
            \textbf N_\textbf B & \textbf N_\mathscr A \textbf W_0^{[k-1]}&\cdots&\textbf N_\mathscr A\textbf W_{n-1}^{[k-1]}\\
        \end{bmatrix}.
    \end{split}
    \end{equation*}
    %
    The rank of $\textbf W$ is bound by the relationship
    \begin{equation*}
    \rank(\begin{bmatrix}
            \textbf D_\mathscr A\!&\!\textbf D_\textbf B
    \end{bmatrix})\!\geq\!    \rank(\begin{bmatrix}\textbf D_\textbf B \!&\! \textbf D_\mathscr A\textbf W_0^{[k-1]} \!&\!\cdots\!&\!\textbf D_\mathscr A\textbf W_{n-1}^{[k-1]}\end{bmatrix}).
    \end{equation*}
    Hence, the upper block of $\textbf W$ has a rank less than $d.$ As a result, $\rank(\textbf W)<n,$ and the nonlinear controllability matrix will also be low rank since it has an identical sparsity structure. Therefore, the pair $(\mathscr A,\textbf B)$ is not structurally controllable. 
    
    
    We have shown the presence of either a dilation or inaccessible node is sufficient to conclude that the system is not structurally controllable. We now show the converse direction. If $(\mathscr A,\textbf B)$ is not structurally controllable, then $\rank(\textbf C) <n.$
    %
    The nonlinear controllability matrix may be low rank when the augmented matrix $\textbf R=\begin{bmatrix}
        \tilde{\textbf A}_{(1)}&\tilde{\textbf B}
    \end{bmatrix}$ has either $\rank(\textbf R)=n$ or $\rank(\textbf R)<n.$
%
    In the case where the augmented matrix $\textbf R$ has a rank less than $n,$ there must exist at least one row of $\textbf R$ that has all zero entries, or $n^{k-1}+m-n+1$ columns that has all zero entries. To satisfy the condition that $\rank(\textbf R)<n$, these rows and columns contain only zero entries, rather than having a nontrivial linear dependence, because if there was a nontrivial linear dependence and there were some nonzero entries, values could be chosen in a fashion consistent with the proof provided by Lee and Markus to increase the rank of $\textbf R.$ Given that these entries are fixed as zeros, the following structures occur in $\Hs(\mathscr A,\textbf B).$ A row of all zeros corresponds to a node that receives no inputs. As a result, the node corresponding to this row is not contained in any hyperedge heads, so there is no walk from the control nodes to this node. Therefore, a nonaccessible node is found. Moreover, at least $n^{k-1}+m-n+1$ columns containing all zeros indicates at most $n-1$ hyperedges in $\Hs(\mathscr A,\textbf B).$ Then, the full set of $n$ nodes of $\Hs(\mathscr A,\textbf B)$ constitutes a dilation since $|\Vs|>|\Es|.$
    It is possible that both a nonaccessible node and dilation exist so that $\rank(\textbf R)<n.$

    In the case where the augmented matrix $\Rv$ has rank $n$ and $\rank(\textbf C)<n$, it will be the case that $\rank(\tilde{\textbf A}_{(1)}\tilde{\textbf B}^{[k-1]})<n,$ since this matrix is a submatrix of the nonlinear controllability matrix.
    %
    Although $\tilde{\textbf A}_{(1)}$ itself may have rank $n$, its image is constrained when acting on inputs from the column space of $\tilde{\textbf B}^{[k-1]}.$ This may occur for two reasons. First, at least one row of $\tilde{\textbf A}_{(1)}\tilde{\textbf B}^{[k-1]}$ may be all zeros. In this case, there are nodes that are nonaccessible on $\Hs(\At,\Bv)$ from a length 1 walk beginning at the control nodes. Second, at least two rows of $\tilde{\textbf A}_{(1)}\tilde{\textbf B}^{[k-1]}$ have a nontrivial linear dependence. If this is the case, there exists a permutation of the coordinates, such that the matrix $\tilde{\textbf A}_{(1)}\tilde{\textbf B}^{[k-1]}$ may be written
        \begin{equation*}
            \tilde{\textbf A}_{(1)}\tilde{\textbf B}^{[k-1]}=\begin{bmatrix}
                \textbf D\\
                \textbf N
            \end{bmatrix},
        \end{equation*}
        where $\textbf D$ denotes the rows that have a nontrivial linear dependence, and $\Nv$ denotes the rows that are linearly independent. Here, $\textbf D$ is a $d\times m^{k}$ matrix, and $\textbf N$ is a $(n-d)\times m^{k}$ matrix, where $\rank(\textbf N)=n-d,$ and $\rank(\textbf D)<d.$ It is possible that $\textbf D$ contains a row of all zeros, in which case a nonaccessible node exist.
        %
        Instead, if all the rows of $\textbf D$ contain a nonzero entry, there must exist at most $d-1$ linearly independent columns of $\textbf D$ to ensure that $\rank(\textbf D)<d.$
        %
        Given that $\textbf D$ is constructed from matrices $(\tilde{\textbf A}_{(1)}, \tilde{\textbf B})$ where the structure is defined but the nonzero entries may take any value, the only way that at most $d-1$ linearly dependent columns are in $\textbf D$ is for $\textbf D$ to contain exactly $d-1$ unique columns, where the columns are unique with respect to which hyperedge tail nodes they are influenced by. As an example, in the tensor of a hypergraph containing a single hyperedge $\{\{v_1\},\{v_2,v_3\}\},$ the tensor structure indicates that the entries $\mathscr A_{123}$ and $\At_{132}$ correspond to the same hyperedge. When unfolded to the $3\times 9$ matrix $\textbf A_{(1)},$ columns 6 and 8 both describe the effects of this hyperedge. 

        
        As the columns of $\textbf D$ indicate nodes contained in hyperedge heads, the set of $d$ nodes whose rows are described by $\textbf D$ has fewer than $d$ hyperedge heads directed at these nodes, so a dilation exists. Recall that $\textbf W_1=\tilde{\textbf A}_{(1)}\tilde{\textbf B}^{[k-1]}$, and it denotes the accessibility of nodes after a one step walk from the control nodes.
    From a walk starting at the sets of nodes indicated by the columns of $\textbf W_j,$ where $\textbf W_0=\Bv,$ the above process illustrates that ensuring $\rank(\textbf W_{j+1})<n$ requires the existence of a dilation or a nonaccessible node.
    Since a nonaccessible node from a walk beginning at $\textbf W_j$ is also nonaccessible from a walk beginning at $\textbf B,$ each step in the walk requires the existence of a dilation of a nonaccessible node.
    %
    Thus, if the system is not structurally controllable then either a dilation or inaccessible node exists.
\end{proof}

\end{document}
